\section{Сравнительный анализ методов развёртывания}
\headingtextsep

Для поставленной задачи можно рассматривать три основных класса методов развёртывания: установщики, контейнерное развёртывание и установку компонентов через пакетные менеджеры. Каждый из этих подходов имеет практическую область применения, однако их пригодность зависит от состава системы, числа сервисов и требований к повторяемости окружения.

\textheadingsep
\subsection{Установщики}
\headingtextsep

Подход с использованием установщиков хорошо подходит для настольных приложений и сравнительно простых серверных систем, поставляемых как один исполняемый модуль. Его сильной стороной является привычный для конечного пользователя формат установки: запуск установочного пакета, копирование файлов, регистрация служб и создание ярлыков.

Для системы «Контур» этот вариант оказывается менее удобным. Приложение состоит из нескольких компонентов, использующих разные технологические стеки и зависящих от отдельной СУБД. В этом случае установщик должен не только развернуть файлы, но и согласованно настроить базу данных, сетевые адреса, порядок запуска процессов и жизненный цикл каждого компонента. Это увеличивает сложность поставки и затрудняет повторяемость развёртывания на разных машинах.

\textheadingsep
\subsection{Пакетные менеджеры}
\headingtextsep

Пакетные менеджеры удобны в ситуациях, когда систему можно разложить на независимые пакеты или когда развёртывается отдельный компонент, например серверный процесс, библиотека или утилита командной строки. Для рассматриваемого проекта такой подход теоретически может быть использован для установки PostgreSQL, Node.js-зависимостей, Rust-бинарников и вспомогательных инструментов.

Однако при использовании пакетных менеджеров возникает проблема сборки общего окружения из разнородных частей. Администратору приходится отдельно устанавливать и согласовывать версии базы данных, backend-сервиса, frontend-приложения и интеграционного моста, а затем вручную настраивать переменные окружения и порядок запуска. В результате повышается вероятность несовпадения версий и появления трудно воспроизводимых ошибок конфигурации.

\textheadingsep
\subsection{Контейнерное развёртывание}
\headingtextsep

Контейнерный подход предполагает упаковку каждого компонента системы вместе с его runtime-зависимостями в отдельный образ. Затем набор образов запускается как согласованная группа контейнеров с фиксированными сетевыми связями, переменными окружения и политиками перезапуска.

Для системы «Контур» этот вариант наиболее естественен. В репозитории уже присутствуют отдельные Dockerfile для API, web-клиента и MCP-моста, а также файл docker-compose.yml, описывающий взаимодействие этих сервисов с контейнером PostgreSQL. Такой способ обеспечивает повторяемость среды, снижает требования к ручной настройке машины и упрощает сопровождение многокомпонентного решения.

\textheadingsep
\subsection{Сравнительный итог}
\headingtextsep

Сравнение показывает, что установщики удобны главным образом для монолитной поставки одного приложения, а пакетные менеджеры --- для установки отдельных компонентов по частям. Для проекта «Контур», представляющего собой многокомпонентную систему из базы данных, backend-сервиса, frontend-клиента и отдельного интеграционного моста, более подходящим является контейнерное развёртывание. Оно лучше остальных способов поддерживает изоляцию зависимостей, воспроизводимость конфигурации и единый порядок запуска всех частей системы.
